\section{Practical design rules from the readings}
\label{sec:practical}

The literature and the experiment can be reduced to a small set of design
rules. I include them because they are easier to apply than a list of
algorithms and they make clear what I would change when working on a real QC
pipeline.

\subsection{Start with the simplest evidence}

A physical range check should run before a complicated model because it is
cheap, deterministic and easy to explain. The same principle applies to
obvious internal contradictions and missing-data checks. More expensive
spatial or forecasting methods are most useful for values that survive the
simple checks. This ordering saves computation and, more importantly, keeps
clearly invalid values from contaminating later references.

\subsection{Use different checks for different fault shapes}

A spike, a frozen sensor and a slow drift are not the same statistical
problem. A step rule is naturally suited to an abrupt jump. A persistence
rule is naturally suited to a loss of movement. A neighbour comparison is
better suited to a sustained offset or drift. The result table supports this
view: the strongest system is the battery because its members fail in
different ways. I would therefore select checks by the faults that matter,
not by which method has the most sophisticated formula.

\subsection{Treat the threshold as part of the model}

It is tempting to describe the algorithm and leave the threshold in a
configuration file. The experiment shows why that is incomplete. Changing
the modified-z threshold from 3.5 to 30 changes the operational behaviour
from complete spike detection with nuisance alarms to silence with missed
spikes. The threshold therefore belongs in the method description together
with the window, sampling rate and reference data used to choose it.

\subsection{Measure reviewer workload, not only accuracy}

A QC system does not end when a flag is produced. Someone may have to inspect
the value, compare other sensors, write a comment or contact a technician.
For that reason I prefer false alarms per day to a percentage that hides the
operational scale. A false-alarm rate that looks small for one series can
become hundreds of flags across a network. A useful evaluation should say
both what faults are found and what attention the system consumes.

\subsection{Keep the original value and the reason for the flag}

The readings by Shafer and Campbell make provenance a central part of the
architecture. An automated check should not silently replace a measurement
with a corrected number or remove it from the archive. The raw value, the QC
verdict, the check that produced it and any later human decision should be
separable. This makes later reprocessing possible when thresholds or methods
change, and it prevents the QC system from hiding scientifically interesting
extremes.

\subsection{Re-evaluate when the context changes}

A threshold tuned for one sampling interval, sensor type or climate should
not automatically be copied to another. The same is true for neighbours and
forecast models. A new site may have different terrain, different natural
variability or different instrument noise. The method can transfer while the
parameter does not. I would therefore treat deployment at a new site as a
small validation exercise rather than as a configuration copy.

These rules are intentionally simple. They do not require a particular
software package or statistical model. They describe how I would make the
system understandable and auditable: simple checks first, several kinds of
evidence, explicit thresholds, measured reviewer cost, preserved provenance,
and local validation.


\subsection*{What I take from each core reading}

The five readings also give five practical lessons that are easy to separate.
Grubbs shows the value of a formally defined statistical decision, but also
how strongly such a decision depends on its assumptions. Iglewicz and
Hoaglin show why robust centre and spread estimates are useful when a few
bad observations can distort the reference. Shafer and colleagues show that
quality control is an operational system made of checks, flags, maintenance
and people, not a single equation. Campbell and colleagues make the
fault-versus-anomaly distinction explicit and argue for preserving raw data
and QC information. Leigh and colleagues show experimentally that method
performance depends on fault type and that combinations are usually more
useful than a single winner.

Together these lessons explain the design used in my experiment: test the
assumptions, make the reference explicit, document the threshold, compare by
fault type, and keep the final verdict interpretable.

% =====================================================================
% 7. Conclusion
% =====================================================================
